\section*{4.3 Operations with Measurable Functions}

Measurable functions can be combined using basic operations, and the resulting function is also measurable. In addition to the usual arithmetic operations, it is also possible to take the limsup and liminf of a sequence of measurable functions. These operations preserve measurability and are useful for constructing new measurable functions from existing ones.

\begin{thmbox}{4.6}
If $f, g : \Omega \to \mathbb{R}$ are $\mathcal{F}$-measurable, then $f + g$, $f - g$, $c \cdot f$, and $f/g$ are measurable, where $c$ is a constant, and in $f/g$, we assume that $g(\omega)$ is nonzero for all $\omega$.
\end{thmbox}

\noindent\textit{Proof} Suppose $f$ and $g$ are measurable functions. Consider
\[
(\Omega, \mathcal{F}) \xrightarrow{\alpha(\omega)} (\mathbb{R}^2, \mathcal{B}(\mathbb{R}^2)) \xrightarrow{\beta(x, y)} (\mathbb{R}, \mathcal{B}(\mathbb{R}))
\]
with $\alpha(\omega)$ defined as $(f(\omega), g(\omega))$, and $\beta(x, y)$ defined as $x + y$. The composite function is $f(\omega) + g(\omega)$. To show that the first function $\alpha$ is measurable, we note that $\mathcal{B}(\mathbb{R}^2)$ can be generated by open rectangles in the form $(a, b) \times (c, d)$ (Theorem 2.8), and
\[
\alpha^{-1}((a, b) \times (c, d)) = f^{-1}((a, b)) \cap g^{-1}((c, d)).
\]
Since both $f^{-1}((a, b))$ and $g^{-1}((c, d))$ are $\mathcal{F}$-measurable (by the assumptions on $f$ and $g$), $\alpha^{-1}((a, b) \times (c, d))$ is $\mathcal{F}$-measurable.

On the other hand, the addition function $\beta(x, y) = x + y$ is Borel measurable, because, for any constant $b$, the pre-image, $\{(x, y) : x + y < b\}$, can be expressed as a countable union of open sets
\[
\bigcup_{\substack{p, q \in \mathbb{Q} \\ p+q < b}} \bigl( \{x < p\} \cap \{y < q\} \bigr).
\]
As a result, the function $\beta(x, y) = x + y$ is measurable by Theorem 4.4. By Theorem 4.3, the composition $\beta \circ \alpha$ is measurable. This proves that $f + g$ is measurable.

Similarly, we can show that the difference, product, and quotient of measurable functions are measurable. To do this, we first show that the linear function $x \mapsto ax$, the square function $x \mapsto x^2$, and the inverse function $x \mapsto 1/x$ are measurable. For any $x$ in $\mathbb{R}$, we check that
\begin{align*}
\{x : ax < b\} &= \begin{cases} \{x : x < b/a\} & \text{if } a > 0, \\ \{x : x > b/a\} & \text{if } a < 0, \end{cases} \\
\{x : x^2 < b\} &= \begin{cases} \{x : -\sqrt{b} < x < \sqrt{b}\} & \text{if } b > 0, \\ \emptyset & \text{otherwise,} \end{cases}
\end{align*}
and
\[
\{x : 1/x < b\} = \begin{cases} \{x : x < 0\} \cup \{x : 1/b < x\} & \text{if } b > 0, \\ \{x : x < 0\} & \text{if } b = 0, \\ \{x : 1/b < x < 0\} & \text{if } b < 0. \end{cases}
\]
Since the sets on the right-hand sides of the displayed equations are all open sets, the sets on the left-hand sides are also open sets. By Theorem 4.4, the linear function, square function, and inverse function are Borel measurable.

We now see that the product function $\beta(x, y) = xy$ is Borel measurable by writing $xy$ as
\[
xy = \frac{1}{2}((x + y)^2 - x^2 - y^2).
\]
Likewise, the quotient function $x/y$ is Borel measurable, because it can be written as
\[
x/y = x(y^{-1}).
\]
We can use a similar proof technique as in the proof that $f + g$ is measurable to show that $f \cdot g$, $f/g$, and $c \cdot f$ are measurable for any measurable functions $f$ and $g$, and any constant $c$. \hfill $\blacksquare$

It is convenient to include the special symbols $\infty$ and $-\infty$ and to work in the extended real number system $\bar{\mathbb{R}} = \mathbb{R} \cup \{\pm\infty\}$. We have the following convention for working with the $\infty$ symbol.

\noindent$\blacktriangleright$ \textbf{Convention Involving Infinity in Probability Theory} The symbols 
\[
\infty - \infty, \ (-\infty) + \infty, \ \infty/\infty, \ \infty/(-\infty), \ (-\infty)/\infty, \text{ and } (-\infty)/(-\infty)
\]
are not defined. However, by convention, we define $0 \cdot \infty$ and $0 \cdot (-\infty)$ to be equal to 0. The other rules for working with infinity follow from common sense. For example,
\[
x \cdot \infty = \begin{cases} \infty & \text{if } x > 0 \\ -\infty & \text{if } x < 0 \\ 0 & \text{if } x = 0. \end{cases}
\]
These conventions are useful when dealing with limits of functions. The following is an example of measurable function that may take $\infty$ as its value.

\textbf{Example 4.3.1 (A Random Variable Whose Value May Be $\infty$)} \\
Suppose $X$, $Y$, and $Z$ are three real-valued random variables defined on a common probability space. Using the machinery we develop so far, such as the continuity of the square root function and the fact that the sum and product of measurable functions are measurable, we can see that $\sqrt{X^2 + Y^2 + Z^2}$ is measurable. The inverse $(X^2 + Y^2 + Z^2)^{-1/2}$ may take $\infty$ as its value, and this happen when $X$, $Y$, and $Z$ are all zero. In this case, we can define $(X^2 + Y^2 + Z^2)^{-1/2}$ as $\infty$.

Recall that the \textit{supremum} of a set $A$ of real numbers, denoted by $\sup A$, is defined as the smallest upper bound of $A$. More precisely, $\sup A$ is the real number $r \in \mathbb{R}$ that satisfies two conditions: (i) $r \ge a$ for all $a \in A$ and (ii) if $u \ge a$ for all $a \in A$, then $u \ge r$. The first condition says that $r$ is an upper bound of $A$, and the second condition says that $r$ is the smallest one. If the set $A$ is unbounded from above, we set $\sup A = \infty$. Thus, $\sup A$ may take a value in the extended real number system. The existence of supremum for any set $A \subseteq \mathbb{R}$ is equivalent to the completeness of the real number system with respect to the absolute value function, and hence, it can be regarded as an axiom of $\mathbb{R}$.

Dually, the \textit{infimum} of a set $A \subseteq \mathbb{R}$, denoted by $\inf A$, is the largest lower bound of $A$. It is the real number $s \in \mathbb{R}$ that satisfies two conditions: (i) $s \le a$ for all $a \in A$, and (ii) if $v \le a$ for all $a \in A$, then $v \le s$. The first condition says that $s$ is a lower bound of $A$, and the second condition says that $s$ is the largest one. When $A$ is not bounded from below, we set $\inf A = -\infty$.

When $A = \{a_1, a_2, a_3, \dots\}$ is a countable set, we use the notation
\[
\sup_i a_i \triangleq \sup\{a_i : i = 1, 2, 3, \dots\}, \text{ and } \inf_i a_i \triangleq \inf\{a_i : i = 1, 2, 3, \dots\}.
\]

\begin{thmbox}{4.7}
Suppose $f_i : \Omega \to \bar{\mathbb{R}}$ is a measurable function, for $i = 1, 2, 3, \dots$. Then the functions $X(\omega) \triangleq \sup_i f_i(\omega)$ and $Y(\omega) \triangleq \inf_i f_i(\omega)$ are measurable.
\end{thmbox}
\noindent\textit{Proof} For any fixed real number $r$, using the definition that the supremum is the smallest upper bound, we obtain
\[
\{\omega : \sup_i f_i(\omega) \le r\} = \bigcap_{i=1}^{\infty} \{\omega : f_i(\omega) \le r\}.
\]
Because the right-hand side is a countable intersection of measurable sets, the left-hand side is also measurable. By Theorem 4.4, the supremum $\sup_i f_i$ is measurable.

We can prove that $\inf_i f_i(\omega)$ is measurable from the relation $\inf_i f_i(\omega) = -\sup_i (-f_i(\omega))$. \hfill $\blacksquare$

Given a sequence of real numbers $(a_i)_{i=1}^{\infty}$, the limsup and liminf of this sequence are defined as
\begin{align*}
\limsup_{i \to \infty} a_i &\triangleq \inf_{j \ge 1} \Big( \sup_{k \ge j} a_k \Big), \\
\liminf_{i \to \infty} a_i &\triangleq \sup_{j \ge 1} \Big( \inf_{k \ge j} a_k \Big).
\end{align*}
In general, the limsup and liminf of a sequence have the following characterizations:
\begin{enumerate}[label=(\alph*)]
    \item $\limsup_k a_k$ is the smallest real number $r$ such that for any $\epsilon > 0$, there is a sufficiently large integer $N$ such that $a_k < r + \epsilon$ for all $k \ge N$.
    \item $\liminf_k a_k$ is the largest real number $s$ such that for any $\epsilon > 0$, there is a sufficiently large integer $N$ such that $a_k > s + \epsilon$ for all $k \ge N$.
\end{enumerate}
For any sequence $(a_i)_{i=1}^{\infty}$, it is generally true that
\[
\liminf_{i \to \infty} a_i \le \limsup_{i \to \infty} a_i.
\]
The inequality becomes an equality when $\lim_{i \to \infty} a_i$ exists.

\textbf{Example 4.3.2 (limsup and liminf)} \\
Consider the numerical sequence $(a_i)_{i=1}^{\infty}$ defined by
\[
a_i \triangleq \begin{cases} 1 + \frac{1}{i} & \text{if } i \text{ is odd,} \\ -1 - \frac{1}{i} & \text{if } i \text{ is even.} \end{cases}
\]
This sequence may be regarded as the interleaving of two sequences that converge to 1 and $-1$, respectively. However, the sequence $(a_i)_{i=1}^{\infty}$ itself does not converge. The limsup and liminf of this sequence are 1 and $-1$, respectively.
\begin{thmbox}{4.8}
Suppose $f_i : \Omega \to \bar{\mathbb{R}}$ is a measurable function, for $i = 1, 2, 3, \dots$. The two functions
\[
X(\omega) \triangleq \limsup_i f_i(\omega) \text{ and } Y(\omega) \triangleq \liminf_i f_i(\omega)
\]
are measurable. Furthermore, if $\lim_{i \to \infty} f_i(\omega)$ exists for all $\omega \in \Omega$, then the function
\[
Z(\omega) = \lim_{i \to \infty} f_i(\omega)
\]
is measurable.
\end{thmbox}

\noindent\textit{Proof} The limsup and liminf of measurable functions are measurable, because
\begin{align*}
\limsup_{i \to \infty} f_i(\omega) &= \inf_{j \ge 1} \sup_{k \ge j} f_k(\omega), \\
\liminf_{i \to \infty} f_i(\omega) &= \sup_{j \ge 1} \inf_{k \ge j} f_k(\omega),
\end{align*}
which can be computed in terms of supremum and infimum. From the previous theorem, we know that the supremum and infimum of measurable functions are measurable.

If the point-wise limit of $f_i$ exists, then the limsup and liminf of $f_i$ coincide. Thus,
\[
\lim_{i \to \infty} f_i(\omega) = \limsup_{i \to \infty} f_i(\omega) = \liminf_{i \to \infty} f_i(\omega)
\]
is measurable. \hfill $\blacksquare$

The operations discussed in this section are sufficient for most applications. For instance, if $X$ and $Y$ are random variables, we can apply the theorems proved in this section to see that
\[
Z = \max(X, Y) + 3 \sin(X) \exp(Y)
\]
is also a random variable. In general, the cdf $F_Z(z) = P(Z \le z)$ may be difficult to compute and may not have a closed-form expression. However, we can at least guarantee that $P(Z \le z)$ is well-defined because we know that $\{\omega : Z(\omega) \le z\}$ is a measurable set.